\section{What the Framework Adds to a Resource Assessment}
\label{sec:implications}

The Philippine example is illustrative; its magnitudes belong to one scenario and one
calibration. What carries over is the kind of information the framework provides. It adds
three things to a resource assessment.

First, it measures the size of each effect on the economy, not only its direction. The
channels can be run alone and together on the same scale, so an assessment can see which
parts of a plan matter for output, consumption and public finances, and how much. In the
example, electricity costs dominate and public investment barely registers; another plan,
with a large publicly financed grid expansion or a carbon tax, would rank the channels
differently.

Second, it follows effects over time. Costs and benefits of a resource plan arrive at
different moments: investment is paid for early, operating savings and cleaner air come
later, and prices may peak during the transition. Because households differ by age, the
framework also shows which generations bear the costs and which receive the benefits.
The effect of a plan in its first decade can differ in size, and even in sign, from its
effect in the long run.

Third, it captures how effects spread through the economy. A change in one sector reaches
others through prices, wages, returns and taxes, and can produce results that a
sector-by-sector calculation would miss. A higher electricity cost that falls mainly on
industry reduces wages and returns throughout the economy; a more capital-intensive
generation technology can reduce the capital held by the sector itself. The economy's
response also feeds back into the plan, changing the demand it must meet and informing
the rate at which it values future costs.

The choice of a resource plan therefore concerns more than the least-cost technology mix. It
also concerns the timing of investment, its financing, and the distribution of costs and
benefits across households and generations. OG--CLEWS brings these questions into the
resource assessment.
